% Generated draft: requires mathematical and source review.
\documentclass[11pt]{article}
\usepackage[margin=25mm]{geometry}
\usepackage{fontspec}
\setmainfont{Latin Modern Roman}
\usepackage{amsmath,amssymb,mathtools,unicode-math}
\setmathfont{Latin Modern Math}
\usepackage{enumitem,needspace,xurl,hyperref}
\setlength{\parindent}{0pt}
\setlength{\parskip}{5pt}
\setlength{\emergencystretch}{3em}
\newcommand{\N}{\mathbb{N}}
\newcommand{\Z}{\mathbb{Z}}
\newcommand{\Q}{\mathbb{Q}}
\newcommand{\R}{\mathbb{R}}
\newcommand{\C}{\mathbb{C}}
\newcommand{\F}{\mathbb{F}}
\newcommand{\A}{\mathbb{A}}
\newcommand{\PP}{\mathbb{P}}
\newcommand{\E}{\mathbb{E}}
\newcommand{\eps}{\varepsilon}
\newcommand{\dd}{\,\mathrm{d}}
\begin{document}
\section*{UnsolvedMath ID: 30003796}
\textbf{Source problem code:} catalogue entry
\subsection{Definitions and mathematical statement}
Let \(K\) be a field, let \(K^{\mathrm{sep}}\) be a separable closure, and write \(\Gamma_K=\operatorname{Gal}(K^{\mathrm{sep}}/K)\). An absolutely almost simple adjoint \(K\)-group is a connected semisimple algebraic \(K\)-group \(G\) with trivial center such that, over an algebraic closure of \(K\), it has no proper connected normal algebraic subgroup of positive dimension. Write \(H^1(K,G)\) for the pointed set of continuous Galois \(1\)-cocycles with values in \(G(K^{\mathrm{sep}})\), modulo the equivalence \(a'_\sigma=g^{-1}a_\sigma\sigma(g)\).

For a set \(V\) of discrete valuations of \(K\), let \(K_v\) denote the completion at \(v\). For each finite subset \(S\subseteq V\), extension of scalars defines
\[
\omega_{S,G}:H^1(K,G)\longrightarrow\prod_{v\in V\setminus S}H^1(K_v,G),\qquad
\alpha\longmapsto(\alpha_{K_v})_{v\in V\setminus S}.
\]
For the set-theoretic formulation below, interpret properness as the requirement that the inverse image of every finite subset is finite, equivalently that every fiber is finite. Thus the required condition is
\[
\text{for every }(\xi_v)_{v\in V\setminus S}\in
\prod_{v\in V\setminus S}H^1(K_v,G),\qquad
\#\{\alpha\in H^1(K,G):\alpha_{K_v}=\xi_v\text{ for every }v\in V\setminus S\}<\infty.
\]
If \(V\setminus S\) is empty, the product is a singleton, so this condition requires \(H^1(K,G)\) itself to be finite. The supplied record does not define properness explicitly; the intended convention requires confirmation.

The supplied clean formulation asks: characterize the fields \(K\) for which there exists a natural set \(V=V(K)\) of discrete valuations such that, for every absolutely almost simple adjoint \(K\)-group \(G\) and every finite \(S\subseteq V\), the map \(\omega_{S,G}\) is proper in the stated sense. In particular, does this property hold for finitely generated fields of good characteristic? Here a finitely generated field means a field finitely generated over its prime subfield.

Two qualifications remain unresolved. First, the record does not specify what makes \(V\) natural or what good characteristic means, including whether the characteristic restriction depends on \(G\). Second, its original statement fixes \(K\) and \(G\) before asking for \(V\), whereas the clean statement requires one choice \(V(K)\) to work for every such \(G\). These quantifier orders need not be equivalent. The displayed target records the clean formulation; the original formulation may instead allow \(V\) to depend on \(G\).
\subsection{Short English statement}
Characterize fields admitting a natural family of discrete valuations for which localization of nonabelian \(H^1\) has finite fibers, even after omitting any finite set of valuations, uniformly over absolutely almost simple adjoint groups. The finitely generated, good-characteristic case is singled out. The meanings of natural and good characteristic, the properness convention, and whether the valuation family may depend on the group require clarification.
\subsection{Sources}
\begin{itemize}
\item Local registry: \texttt{data/unsolvedmath/problems.json}, supplied record with permanent UnsolvedMath ID \(30003796\), exact source problem code \texttt{catalogue entry}, title \emph{Properness of Localization Maps for Algebraic Groups}, in the Oberwolfach Reports Open Problems collection. No individual proposer is identified; the supplied proposed year is 2018.
\item Supplied report citation: \emph{Quadratic Forms and Related Structures over Fields (2019)}, Oberwolfach Reports, pp.~1235--1291, DOI \texttt{10.4171/owr/2018/21}. \url{https://doi.org/10.4171/owr/2018/21}. The supplied background describes these reports as records of mathematical workshops and their open-problem discussions.
\item Supplied literature source, labeled ``spinor groups with good reduction'': \url{https://arxiv.org/abs/1707.08062}.
\item Supplied literature source, labeled ``good-reduction survey'': \url{https://arxiv.org/abs/2005.05484}.
\item Supplied literature source, labeled ``workshop report'': \url{https://ems.press/content/serial-article-files/46742}.
\item Literature assessment checked at \texttt{2026-08-21}: the supplied status is \texttt{partially\_solved}. The supplied assessment states that the properness question is settled for number fields and several group classes, with partial higher-dimensional results, but not for arbitrary fields or all finitely generated good-characteristic groups. This is supplied dated curation derived from the OWR record, not an independently verified literature assessment; its status and cited claims require independent verification.
\item The supplied statement status is \texttt{corrected\_verified}; its verification note says that the clean statement was checked against the official Oberwolfach report and preserves its hypotheses, notation, and open target. This is a supplied verification claim, not verification performed for this draft. The report citation gives 2019, while the proposed year and year tag give 2018 and the DOI contains 2018; these supplied date distinctions are retained without resolving their bibliographic meaning.
\end{itemize}
\end{document}
